\documentclass[12pt]{letter}
% Letter form, formatted for CSUC letterhead 

\begin{document}

\signature{Dr. Eric Ayars \\ Professor of Physics \\ CSU Chico Department of Physics}

\begin{letter}{}
	
\opening{Dear Dr. Essick,}

Please convey our gratitude to both reviewers for their time and effort in reviewing our manuscript.
We have carefully considered their comments and have made additions to the manuscript accordingly.

The most significant change is the increased emphasis on how this device can be used in labs beyond the first year. 
The reviewers' comments made clear to me that I was not selling it as well as I could have: things that were obvious to me needed to be brought out in the text.

We explicitly pointed out that the device can report exact position at any time during the drive cycle. 
This capability allows for the capture of drive phase information, and we called out that phase information in figure 4.
We suggested specifically upper-divison experiments that can be done with this device, such as using the drive-position data to generate a complete set of Poincar\'e sections for chaotic systems.
We also suggested the possibility of using this hardware with alternate firmware to generate non-sinusoidal outputs, or multi-frequency outputs for use in Advanced Lab experiments, or using further modifications to the hardware to allow experimental feedback to the controller. 
All these are definitely Advanced Lab applications, which should address the concerns of Reviewer 3 (and Beth) and put the article firmly into the scope of AJP.

I've dug into the suggestion of Reviewer 1 that we should include something showing the improved response of a simple resonance system to this improved driver. 
The problem I've run into is that the simple cart-spring system we've been testing it on is not particularly sensitive to the drive harmonics! 
I did a frequency sweep from $f_o/5$ to $2f_o$, and as expected there is a cart response at $f_o/2$ where the $2f_o$ harmonic causes a slight resonance bump, but the bump is nearly buried in the noise and not particularly compelling. 
For this system, it's not much. 
We've been running with the apparatus, though: I currently have a  summer student building a chaotic system driven by this device, and that system would be much more sensitive to the drive harmonics.
But we don't have that data yet; it'll be another paper entirely!
I am hoping that we can leave that bit as is, with the mention of the measured 20dB reduction in that first overtone.

I am also including a .zip file containing the SolidWorks and .STL files for the device, a schematic, and the firmware code for the microcontroller.

\closing{Sincerely,}

\end{letter}

\end{document}
		

